%% ============================================================
%%  Abstract & Keywords — NMI framing, IEEEtran syntax
%%  Content mirrors 006_abstract_nmi.tex (Nature build) but uses
%%  \begin{abstract} / \begin{IEEEkeywords} so it compiles under
%%  the IEEEtran class (main.tex). Keep the two files in sync.
%% ============================================================

\begin{abstract}
As AI models scale beyond fast-memory capacity, the performance of
large-scale inference is limited not by computation but by how intelligent
systems coordinate data across a hierarchical memory system. Here we show
that this coordination is not continuously optimisable but is governed by
intrinsic, regime-dependent limits. Characterising any inference system by
two dimensionless ratios---the fast-memory residency ratio ($R_C$) and the
transfer-pressure ratio ($R_B$)---we identify three operational
regimes---capacity-limited, coordination-dominated, and I/O-limited---separated
by abrupt, phase-like transitions. Across these boundaries the ranking of
optimisation strategies inverts: a method that lowers latency by 24\% in one
regime raises it by 8--12\% in another. We derive structural lower bounds
proving the transitions are inevitable, calibrate the boundaries across five
hardware platforms, and confirm the principle on language, vision, and
retrieval-augmented workloads. These results establish regime-dependent
behaviour as a general principle of memory-bound machine intelligence,
reframing AI-inference scaling from continuous tuning to regime-aware design.
\end{abstract}

\begin{IEEEkeywords}
Machine intelligence at scale, AI inference, memory-bound computation,
regime transitions, phase-like behaviour, Large Language Models,
hierarchical memory systems
\end{IEEEkeywords}
